
\chapter{连续系统振动——弦与杆}

\intro{本章从离散系统的局限出发，引入连续系统的概念。你将学习用偏微分方程描述弦的横向振动、杆的纵向振动和扭转振动，掌握分离变量法求解固有频率和振型，并理解模态叠加法在连续系统中的推广。}

%%%%%%%%%%%%%%%%%%%%%%%%%%%%%%%%%%%%%%%%%%%%%%%%%%%%%%%%
\section{连续系统与离散系统的区别}

\subsection{从离散到连续——自由度概念的扩展}

前六章讨论的离散系统（单自由度、多自由度）由有限个集中质量、弹簧和阻尼器组成，其运动由有限个常微分方程（ODE）描述。然而，工程中的许多结构——如弦、杆、梁、板和壳——其质量和弹性是连续分布的。这类\textbf{连续系统}（continuous system）拥有\textbf{无限多个自由度}，其运动必须用\textbf{偏微分方程}（PDE）来描述。

两者的本质区别：

\begin{mydefinition}{离散系统 vs 连续系统}
\begin{itemize}
\item \textbf{离散系统}：有限个常微分方程;运动由有限个广义坐标 $q_1(t),\dots,q_n(t)$ 描述;固有频率有有限个。
\item \textbf{连续系统}：偏微分方程;运动由空间和时间的函数 $y(x,t)$ 描述;固有频率有无穷多个。
\end{itemize}
\end{mydefinition}

从数学上看，离散系统求解的是代数特征值问题 $\det(\mathbf{K} - \omega^2\mathbf{M})=0$;而连续系统求解的是微分算子的特征值问题（Sturm-Liouville 问题），涉及边界条件。

\begin{mydefinition}{连续系统的建模思路}
连续系统的分析遵循以下步骤：(1) 取微元体，建立力平衡方程或能量关系;(2) 得到关于位移 $y(x,t)$ 的 PDE;(3) 施加边界条件和初始条件;(4) 用分离变量法或积分变换法求解。
\end{mydefinition}

\subsection{边界条件的核心作用}

离散系统的固有频率完全由系统内部参数（$m, k, c$）决定。而在连续系统中，\textbf{边界条件}是解的内在组成部分——相同的 PDE 在不同边界条件下会给出完全不同的固有频率和振型。例如，两端固定的弦与一端固定一端自由的弦，其固有频率和振型完全不同。这是连续系统与离散系统最显著的区别之一。

--------------------------------------------------------------------
\section{弦的横向振动}

\subsection{物理模型与基本假设}

考虑一根张紧的柔软弦，长度为 $L$，单位长度质量为 $\rho$，在张力 $T$ 作用下处于平衡位置。弦做微小横向振动，位移记为 $y(x,t)$。基本假设：

\begin{enumerate}
    \item \textbf{小挠度假设}：位移 $y$ 远小于弦长 $L$，斜率 $\partial y/\partial x$ 很小;
    \item \textbf{均匀性假设}：$\rho$ 和 $T$ 沿弦长不变;
    \item \textbf{柔性假设}：弦不抵抗弯曲，张力沿弦的方向;
    \item \textbf{无阻尼假设}：忽略能量耗散。
\end{enumerate}

\subsection{波动方程的推导}

在 $x$ 处取长度为 $dx$ 的微段，如图 \ref{fig:string-element} 所示。微段两端受张力 $T$ 作用，方向沿弦的切线方向。在微小振动的条件下，$\sin\theta \approx \tan\theta = \partial y/\partial x$。

\begin{figure}[H]\centering
\begin{tikzpicture}[vib_style/.style={thick},vib_arrow/.style={-Stealth,thick}]
\draw[vib_style] (0,0) -- (10,0) node[right] {$x$};
\draw[vib_style] (0,0) -- (0,4) node[above] {$y$};
\draw[vib_style] (2,0) -- (2,-0.2) node[below] {$x$};
\draw[vib_style] (6,0) -- (6,-0.2) node[below] {$x+dx$};
\draw[vib_style,dashed] (2,0) -- (2,2.5);
\draw[vib_style,dashed] (6,0) -- (6,1.8);
\draw[blue,ultra thick] plot[smooth] coordinates {(1.5,2.3) (2,2.5) (3,2.2) (4,1.5) (5,1.2) (6,1.8) (6.5,2.0)};
\draw[vib_arrow,red] (2,2.5) -- (1.3,2.9);
\draw[vib_arrow,red] (6,1.8) -- (7.0,1.4);
\node at (1.3,3.1) {$T$};
\node at (7.0,1.2) {$T$};
\draw (2.5,2.6) node {$\theta$};
\draw (5.5,1.5) node {$\theta + d\theta$};
\draw[|-|] (1.8,3.0) -- (6.2,3.0) node[midway,above] {$dx$};
\end{tikzpicture}
\caption{弦微段的受力分析。两端张力方向沿切线方向}
\label{fig:string-element}
\end{figure}

微段在 $y$ 方向的受力为：
\[
F_y = T\sin(\theta + d\theta) - T\sin\theta \approx T\frac{\partial}{\partial x}\left(\frac{\partial y}{\partial x}\right)dx = T\frac{\partial^2 y}{\partial x^2}dx
\]

由牛顿第二定律 $F = ma$，微段质量 $\rho dx$，加速度 $\partial^2 y/\partial t^2$：
\[
\rho dx \cdot \frac{\partial^2 y}{\partial t^2} = T\frac{\partial^2 y}{\partial x^2}dx
\]

约去 $dx$ 得到经典的\textbf{一维波动方程}：

\begin{myformula}
\[
\frac{\partial^2 y}{\partial t^2} = c^2\frac{\partial^2 y}{\partial x^2}, \quad c = \sqrt{\frac{T}{\rho}}
\]
\end{myformula}

其中 $c$ 为波速（wave speed），量纲为 m/s。张力越大，波速越高;线密度越大，波速越低。

\begin{remark}
波动方程是双曲型 PDE 的典范。它描述了波的传播：信息以有限速度 $c$ 在弦上传播，而非瞬时传递。这是连续介质力学与刚体力学的重要区别。
\end{remark}

\subsection{分离变量法}

假设解可以写为空间函数与时间函数的乘积：
\[
y(x,t) = Y(x) \cdot T(t)
\]

代入波动方程：
\[
Y(x)\ddot{T}(t) = c^2 Y''(x) T(t)
\]

两边除以 $c^2 Y(x) T(t)$（当 $Y\neq0, T\neq0$ 时）：
\[
\frac{\ddot{T}(t)}{c^2 T(t)} = \frac{Y''(x)}{Y(x)} = -\beta^2 \quad (\text{常数})
\]

这里引入分离常数 $-\beta^2$ 的原因：等号左边仅是 $t$ 的函数，右边仅是 $x$ 的函数，两者相等意味着必须等于同一个常数。取 $-\beta^2$ 保证了时间方程的解是简谐振动（而非指数增长），符合物理实际。

\subsection{空间方程与时间方程}

由此得到两个常微分方程：

\begin{myformula}
\[\begin{aligned}
\text{空间方程：}\quad & Y''(x) + \beta^2 Y(x) = 0 \\
\text{时间方程：}\quad & \ddot{T}(t) + \omega^2 T(t) = 0, \quad \omega = \beta c
\end{aligned}\]
\end{myformula}

空间方程的通解为：
\[
Y(x) = A\sin(\beta x) + B\cos(\beta x)
\]

时间方程的通解为：
\[
T(t) = C\sin(\omega t) + D\cos(\omega t)
\]

其中 $A,B,C,D$ 为由边界条件和初始条件确定的积分常数。$\omega = \beta c$ 为\textbf{固有角频率}，$\beta$ 为\textbf{波数}。

\subsection{两端固定的弦}

这是最经典的情形——如吉他弦、小提琴弦。

\textbf{边界条件}：$y(0,t)=0$（左端固定），$y(L,t)=0$（右端固定），即 $Y(0)=0$，$Y(L)=0$。

代入 $Y(0)=0$：
\[
Y(0) = A\cdot0 + B\cdot1 = B = 0 \quad\Rightarrow\quad B = 0
\]

因此 $Y(x) = A\sin(\beta x)$。

代入 $Y(L)=0$：
\[
Y(L) = A\sin(\beta L) = 0, \quad A \neq 0 \quad\Rightarrow\quad \sin(\beta L) = 0
\]

这给出\textbf{频率方程}（frequency equation）：
\[
\beta_n L = n\pi, \quad n = 1,2,3,\dots
\]

因此固有波数和固有频率为：

\begin{myformula}
\[
\beta_n = \frac{n\pi}{L}, \quad \omega_n = \beta_n c = \frac{n\pi}{L}\sqrt{\frac{T}{\rho}}, \quad n=1,2,3,\dots
\]
\end{myformula}

固有振型（mode shape）为：
\[
Y_n(x) = \sin\left(\frac{n\pi x}{L}\right), \quad n=1,2,3,\dots
\]

\begin{figure}[H]\centering
\begin{tikzpicture}[scale=1.4,vib_style/.style={thick},vib_arrow/.style={-Stealth,thick}]
\draw[vib_style] (0,0) -- (4,0) node[right] {$x$};
\draw[vib_style] (0,0) -- (0,1.5) node[above] {$y$};
\draw[vib_style,fill=black] (0,0) circle (0.06);
\draw[vib_style,fill=black] (4,0) circle (0.06);
\node at (0,-0.2) {$0$};
\node at (4,-0.2) {$L$};

\draw[domain=0:4,smooth,variable=\x,blue,ultra thick] plot ({\x},{sin(pi*\x/4 r)});
\node[blue] at (4.5,1.0) {$n=1$};

\draw[domain=0:4,smooth,variable=\x,red,ultra thick] plot ({\x},{0.9*sin(2*pi*\x/4 r)});
\node[red] at (4.5,0.9) {$n=2$};

\draw[domain=0:4,smooth,variable=\x,green!60!black,ultra thick] plot ({\x},{0.8*sin(3*pi*\x/4 r)});
\node[green!60!black] at (4.5,0.8) {$n=3$};

\node at (2,-0.6) {两端固定弦的前三阶振型};
\end{tikzpicture}
\caption{两端固定弦的前三阶固有振型。$n=1$ 为基频振型，$n=2,3$ 为高阶振型。振型在两端恒为零}
\label{fig:string-modes-fixed-fixed}
\end{figure}

物理意义：$n=1$ 给出\textbf{基频}（fundamental frequency），弦以半个波长振动;$n=2$ 为第一泛音（octave），弦以完整波长振动;$n$ 每增加 1，固有频率增加一倍（谐波关系）。这使得弦乐器产生悦耳的和声——因为所有泛音都是基频的整数倍。

\subsection{一端固定一端自由的弦}

\textbf{边界条件}：$y(0,t)=0$（固定端），$T \cdot \partial y/\partial x\big|_{x=L}=0$ 即 $Y'(L)=0$（自由端不受横向力）。

由 $Y(0)=0$ 得 $B=0$，即 $Y(x)=A\sin(\beta x)$。

由 $Y'(L)=0$：
\[
Y'(L) = A\beta\cos(\beta L) = 0 \quad\Rightarrow\quad \cos(\beta L)=0
\]

频率方程：$\beta_n L = (2n-1)\pi/2$，$n=1,2,3,\dots$

固有频率和振型：
\[
\omega_n = \frac{(2n-1)\pi}{2L}\sqrt{\frac{T}{\rho}}, \quad Y_n(x)=\sin\left(\frac{(2n-1)\pi x}{2L}\right)
\]

\begin{remark}
与两端固定不同，固定-自由边界条件下固有频率不是基频的整数倍（出现半奇数倍），振型在自由端取到极值而非节点。
\end{remark}

\begin{figure}[H]\centering
\begin{tikzpicture}[scale=1.4,vib_style/.style={thick},vib_arrow/.style={-Stealth,thick}]
\draw[vib_style] (0,0) -- (4,0) node[right] {$x$};
\draw[vib_style] (0,0) -- (0,1.5) node[above] {$y$};
\draw[vib_style,fill=black] (0,0) circle (0.06);
\draw[vib_style] (4,0) circle (0.06);
\node at (0,-0.2) {$0$};
\node at (4,-0.2) {$L$};

\draw[domain=0:4,smooth,variable=\x,blue,ultra thick] plot ({\x},{sin(pi*\x/8 r)});
\node[blue] at (4.5,0.5) {$n=1$};

\draw[domain=0:4,smooth,variable=\x,red,ultra thick] plot ({\x},{0.9*sin(3*pi*\x/8 r)});
\node[red] at (4.5,0.4) {$n=2$};

\draw[domain=0:4,smooth,variable=\x,green!60!black,ultra thick] plot ({\x},{0.8*sin(5*pi*\x/8 r)});
\node[green!60!black] at (4.5,0.3) {$n=3$};

\node at (2,-0.6) {一端固定一端自由弦的前三阶振型};
\end{tikzpicture}
\caption{固定-自由弦的前三阶固有振型。自由端斜率为零，振型取到极值}
\label{fig:string-modes-fixed-free}
\end{figure}

\subsection{两端自由的弦}

\textbf{边界条件}：$Y'(0)=0$，$Y'(L)=0$。

由 $Y'(0)=0$ 得 $A=0$，即 $Y(x)=B\cos(\beta x)$。由 $Y'(L)=0$：
\[
-B\beta\sin(\beta L)=0 \quad\Rightarrow\quad \sin(\beta L)=0
\]

频率方程与两端固定的弦相同：$\beta_n L = n\pi$。但基频对应 $n=0$（$\beta_0=0$），此时 $\omega_0=0$，振型 $Y_0(x)=\text{const}$——对应\textbf{刚体平移模态}（rigid-body mode）。

固有振型：$Y_n(x)=\cos(n\pi x/L)$，$n=0,1,2,\dots$。$n=0$ 为刚体模态，$n\ge 1$ 为弹性模态。

\begin{remark}
任何具有自由-自由边界的系统都存在零频率的刚体模态。这是连续系统中需要特别注意的现象——两端自由的梁、浮动的板都有类似特性。
\end{remark}

\subsection{振型的正交性}

固有振型函数 $Y_n(x)$ 具有重要的\textbf{正交性质}（orthogonality），这在模态叠加法中至关重要。

\begin{theorem}[弦振型的正交性]
不同阶次的固有振型关于质量分布加权正交：
\[
\int_0^L Y_m(x) Y_n(x)\,dx = 0, \quad m \neq n
\]
\end{theorem}

\begin{proof}
$Y_m$ 和 $Y_n$ 分别满足空间方程：
\[
Y_m'' + \beta_m^2 Y_m = 0, \quad Y_n'' + \beta_n^2 Y_n = 0
\]
将第一式乘 $Y_n$，第二式乘 $Y_m$，相减并在 $[0,L]$ 上积分：
\[
\int_0^L (Y_m'' Y_n - Y_n'' Y_m)\,dx + (\beta_m^2-\beta_n^2)\int_0^L Y_m Y_n\,dx = 0
\]
第一项分部积分两次：
\[
\int_0^L (Y_m'' Y_n - Y_n'' Y_m)\,dx = [Y_m'Y_n - Y_n'Y_m]_0^L
\]
对于固定端（$Y=0$）或自由端（$Y'=0$）等齐次边界条件，上式边界项为零。由于 $\beta_m\neq\beta_n$（$m\neq n$），必有
$\int_0^L Y_m Y_n\,dx = 0$。
\end{proof}

同样可以证明关于刚度的正交性：
\[
\int_0^L Y_m'(x) Y_n'(x)\,dx = 0, \quad m \neq n
\]

振型的归一化：通常取 $\int_0^L Y_n^2(x)\,dx = L/2$（对两端固定的弦）。

--------------------------------------------------------------------
\section{杆的纵向振动}

\subsection{物理模型与基本假设}

考虑一均匀弹性杆，长度 $L$，截面积 $A$，密度 $\rho$，弹性模量 $E$。杆做纵向（轴向）振动，位移记为 $u(x,t)$。

基本假设：(1) 截面在振动中保持平面（平截面假设）;(2) 杆细长，横向惯性效应可忽略;(3) 应力-应变关系服从 Hooke 定律 $\sigma = E\varepsilon$。

\subsection{波动方程的推导}

取 $x$ 处的微段 $dx$，截面受力如图 \ref{fig:rod-element}。

\begin{figure}[H]\centering
\begin{tikzpicture}[vib_style/.style={thick},vib_arrow/.style={-Stealth,thick}]
\draw[vib_style] (0,0) -- (8,0) node[right] {$x$};
\draw[vib_style,fill=gray!20] (2,-1.5) rectangle (2.4,1.5);
\draw[vib_style,fill=gray!20] (5.6,-1.5) rectangle (6,1.5);
\draw[dashed] (2.2,-1.5) -- (2.2,1.5);
\draw[dashed] (5.8,-1.5) -- (5.8,1.5);
\draw[vib_arrow,red] (1.5,0) -- (2.2,0) node[midway,above] {$P$};
\draw[vib_arrow,red] (6.5,0) -- (5.8,0) node[midway,above] {$P+dP$};
\draw (2.2,-0.3) -- (2.2,-0.2) node[below] {$x$};
\draw (5.8,-0.3) -- (5.8,-0.2) node[below] {$x+dx$};
\end{tikzpicture}
\caption{杆微段纵向受力。左端受 $P$，右端受 $P+dP$}
\label{fig:rod-element}
\end{figure}

微段的应变：$\varepsilon = \partial u/\partial x$。由 Hooke 定律：$\sigma = E\varepsilon = E\partial u/\partial x$。截面力：$P = \sigma A = EA \partial u/\partial x$。

牛顿第二定律应用于微段：
\[
(P+dP) - P = \rho A dx \cdot \frac{\partial^2 u}{\partial t^2}
\]
即 $dP = (\partial P/\partial x)dx = EA (\partial^2 u/\partial x^2)dx$，代入得：

\begin{myformula}
\[
\frac{\partial^2 u}{\partial t^2} = c_L^2\frac{\partial^2 u}{\partial x^2}, \quad c_L = \sqrt{\frac{E}{\rho}}
\]
\end{myformula}

$c_L$ 为杆中的纵波波速。例如钢材 $E\approx 210$ GPa，$\rho\approx 7850$ kg/m$^3$，$c_L\approx 5172$ m/s。

\begin{remark}
杆的纵向振动方程在形式上与弦的横向振动方程完全相同——都是一维波动方程。区别在于波速的物理来源：弦的波速源于张力与线密度之比，杆的波速源于弹性模量与体密度之比。
\end{remark}

\subsection{不同边界条件下的解}

\subsubsection{固定-固定杆}

边界条件：$u(0,t)=0$，$u(L,t)=0$。

固有频率和振型：
\[
\omega_n = \frac{n\pi}{L}\sqrt{\frac{E}{\rho}}, \quad U_n(x)=\sin\left(\frac{n\pi x}{L}\right), \quad n=1,2,3,\dots
\]

\begin{figure}[H]\centering
\begin{tikzpicture}[scale=1.3,vib_style/.style={thick},vib_arrow/.style={-Stealth,thick}]
\draw[vib_style] (0,-1.5) -- (0,1.8) node[above] {$u$};
\draw[vib_style] (0,0) -- (5,0) node[right] {$x$};
\draw[vib_style,fill=black] (0,0) circle (0.06);
\draw[vib_style,fill=black] (5,0) circle (0.06);
\node at (0,-0.2) {$0$};
\node at (5,-0.2) {$L$};

\draw[domain=0:5,smooth,variable=\x,blue,ultra thick] plot ({\x},{1.2*sin(pi*\x/5 r)});
\node at (5.5,1.2) [blue] {$n=1$};
\draw[domain=0:5,smooth,variable=\x,red,ultra thick] plot ({\x},{1.0*sin(2*pi*\x/5 r)});
\node at (5.5,1.0) [red] {$n=2$};
\draw[domain=0:5,smooth,variable=\x,green!60!black,ultra thick] plot ({\x},{0.8*sin(3*pi*\x/5 r)});
\node at (5.5,0.8) [green!60!black] {$n=3$};

\node at (2.5,-1.0) {固定-固定杆的纵向振型};
\end{tikzpicture}
\caption{固定-固定杆的前三阶纵向振型。$n=1$ 时中点振幅最大，$n=2$ 时在 $x=L/2$ 处出现节点}
\label{fig:rod-modes-fixed-fixed}
\end{figure}

\subsubsection{固定-自由杆}

边界条件：$u(0,t)=0$（固定端位移为零），$EA \cdot \partial u/\partial x\big|_{x=L}=0$（自由端正应力为零），即 $U(0)=0$，$U'(L)=0$。

固有频率和振型：
\[
\omega_n = \frac{(2n-1)\pi}{2L}\sqrt{\frac{E}{\rho}}, \quad U_n(x)=\sin\left(\frac{(2n-1)\pi x}{2L}\right), \quad n=1,2,3,\dots
\]

\subsubsection{自由-自由杆}

边界条件：$U'(0)=0$，$U'(L)=0$。

固有频率和振型：
\[
\omega_n = \frac{n\pi}{L}\sqrt{\frac{E}{\rho}}, \quad U_n(x)=\cos\left(\frac{n\pi x}{L}\right), \quad n=0,1,2,\dots
\]

其中 $n=0$ 为刚体模态（$\omega_0=0$，$U_0(x)=1$）。

\begin{figure}[H]\centering
\begin{tikzpicture}[scale=1.3,vib_style/.style={thick},vib_arrow/.style={-Stealth,thick}]
\draw[vib_style] (0,-1.8) -- (0,1.8) node[above] {$u$};
\draw[vib_style] (0,0) -- (5,0) node[right] {$x$};
\draw[vib_style] (0,0) circle (0.06);
\draw[vib_style] (5,0) circle (0.06);
\node at (0,-0.2) {$0$};
\node at (5,-0.2) {$L$};

\draw[domain=0:5,smooth,variable=\x,blue,ultra thick] plot ({\x},{1.2*cos(pi*\x/5 r)});
\node at (5.5,1.2) [blue] {$n=1$};
\draw[domain=0:5,smooth,variable=\x,red,ultra thick] plot ({\x},{1.0*cos(2*pi*\x/5 r)});
\node at (5.5,1.0) [red] {$n=2$};
\draw[domain=0:5,smooth,variable=\x,green!60!black,ultra thick] plot ({\x},{0.8*cos(3*pi*\x/5 r)});
\node at (5.5,0.8) [green!60!black] {$n=3$};

\node at (2.5,-1.3) {自由-自由杆的纵向振型（不含刚体模态）};
\end{tikzpicture}
\caption{自由-自由杆前三阶弹性纵向振型。两端位移不为零，斜率为零}
\label{fig:rod-modes-free-free}
\end{figure}

\subsection{边界条件对固有频率的影响总结}

四种典型边界条件及其固有频率：

\begin{table}[H]\centering
\caption{不同边界条件下杆的纵向振动固有频率}
\label{tab:rod-frequencies}
\begin{tabular}{ccc}
\toprule
边界条件 & 频率方程 & 固有频率 $\omega_n$ \\
\midrule
固定-固定 & $\sin(\beta L)=0$ & $\frac{n\pi}{L}\sqrt{\frac{E}{\rho}}$ \\
固定-自由 & $\cos(\beta L)=0$ & $\frac{(2n-1)\pi}{2L}\sqrt{\frac{E}{\rho}}$ \\
自由-自由 & $\sin(\beta L)=0$ & $\frac{n\pi}{L}\sqrt{\frac{E}{\rho}}$（$n\ge0$）\\
\bottomrule
\end{tabular}
\end{table}

--------------------------------------------------------------------
\section{杆的扭转振动}

\subsection{物理模型与基本假设}

考虑一均匀圆截面杆做扭转振动，扭转角记为 $\theta(x,t)$。基本假设：(1) 圆截面在扭转后仍保持平面（平截面假设）;(2) 截面绕轴线刚性旋转，无翘曲;(3) 材料的剪切模量为 $G$。

\subsection{扭转波动方程}

微段 $dx$ 的扭转惯性力矩为 $\rho I_p dx \cdot \partial^2\theta/\partial t^2$，其中 $I_p$ 为截面极惯性矩。弹性恢复力矩为 $GI_p \partial\theta/\partial x$。两端力矩之差给出：

\begin{myformula}
\[
\frac{\partial^2\theta}{\partial t^2} = c_T^2\frac{\partial^2\theta}{\partial x^2}, \quad c_T = \sqrt{\frac{G}{\rho}}
\]
\end{myformula}

$c_T$ 为扭转波速（剪切波速）。注意到对大多数金属材料，$G \approx E/[2(1+\nu)]$（$\nu\approx 0.3$），因 此 $c_T < c_L$。

扭转振动方程在数学形式上与纵向振动完全相同。因此，相同边界条件下扭转振动的固有频率公式与纵向振动一致，只需将 $c_L$ 替换为 $c_T$。

\begin{remark}
对比三种波动：
\begin{itemize}
    \item \textbf{弦的横波}：$c=\sqrt{T/\rho}$（张力驱动）
    \item \textbf{杆的纵波}：$c_L=\sqrt{E/\rho}$（弹性驱动）
    \item \textbf{杆的扭转波}：$c_T=\sqrt{G/\rho}$（剪切驱动）
\end{itemize}
三者的控制方程都是 $\partial^2\phi/\partial t^2 = c^2\partial^2\phi/\partial x^2$——这说明一维波动方程是连续介质力学中波传播的普适形式。
\end{remark}

--------------------------------------------------------------------
\section{模态叠加法用于连续系统}

\subsection{模态展开}

与多自由度系统中的模态叠加法类似，连续系统的任意响应可以展开为固有振型的无穷级数：

\begin{myformula}
\[
y(x,t) = \sum_{n=1}^{\infty} Y_n(x)\, q_n(t)
\]
\end{myformula}

其中 $Y_n(x)$ 为归一化的固有振型，$q_n(t)$ 为第 $n$ 阶模态坐标（modal coordinate）。

将展开式代入弦的受迫振动方程：
\[
\rho\frac{\partial^2 y}{\partial t^2} - T\frac{\partial^2 y}{\partial x^2} = f(x,t)
\]
两边乘以 $Y_m(x)$，在 $[0,L]$ 上积分，利用正交性 $\int_0^L \rho Y_m Y_n dx = M_n\delta_{mn}$ 和 $\int_0^L T Y_m' Y_n' dx = K_n\delta_{mn}$，得到解耦的模态方程：

\begin{myformula}
\[
M_n \ddot{q}_n(t) + K_n q_n(t) = Q_n(t), \quad n=1,2,3,\dots
\]
\end{myformula}

其中：
\begin{itemize}
    \item 模态质量：$M_n = \int_0^L \rho Y_n^2(x)\,dx$
    \item 模态刚度：$K_n = \int_0^L T [Y_n'(x)]^2\,dx$（或 $EA [U_n'(x)]^2$）
    \item 模态力：$Q_n(t) = \int_0^L f(x,t)Y_n(x)\,dx$
    \item 固有频率：$\omega_n = \sqrt{K_n/M_n}$
\end{itemize}

\begin{remark}
模态叠加法的威力在于：它将无限自由度的 PDE 转化为无穷多个解耦的 ODE，每个 ODE 恰好是单自由度简谐振子的方程。在工程实践中，取前若干阶模态即可获得满意的精度。
\end{remark}

\subsection{含阻尼的受迫振动}

当系统包含粘性阻尼时，模态方程变为：
\[
M_n \ddot{q}_n + 2\zeta_n\omega_n M_n \dot{q}_n + K_n q_n = Q_n(t)
\]

对于零初始条件的简谐激励 $f(x,t)=F(x)\sin(\Omega t)$，稳态响应为：
\[
y(x,t) = \sum_{n=1}^{\infty} \frac{Y_n(x)}{K_n} \cdot \frac{1}{\sqrt{(1-\Omega^2/\omega_n^2)^2 + (2\zeta_n\Omega/\omega_n)^2}} \cdot \sin(\Omega t - \phi_n)
\]

这表明连续系统在与某阶固有频率相近的频率下会发生\textbf{共振}——该阶模态的响应占主导地位。

\subsection{集中力作用下的响应}

若在 $x=x_0$ 处作用集中力 $F(t)=F_0\sin(\Omega t)$，模态力为：
\[
Q_n(t) = \int_0^L F_0\sin(\Omega t)\delta(x-x_0)Y_n(x)\,dx = F_0 Y_n(x_0)\sin(\Omega t)
\]

稳态响应：
\[
y(x,t) = F_0 \sum_{n=1}^{\infty} \frac{Y_n(x)Y_n(x_0)}{K_n(1-\Omega^2/\omega_n^2)}\sin(\Omega t)
\]

注意：若激振点位于某阶振型的节点（$Y_n(x_0)=0$），则该阶模态不被激发——这是模态测试中布置传感器的关键依据。

--------------------------------------------------------------------
\section{数值模拟与代码实现}

\subsection{弦振动的有限差分解}

弦振动的波动方程可用中心差分法离散。将空间和时间分别离散为 $N_x$ 和 $N_t$ 个点：

\begin{mycode}{弦振动的有限差分法（Python）}
\begin{lstlisting}[language=Python]
import numpy as np
import matplotlib.pyplot as plt
from matplotlib.animation import FuncAnimation
import matplotlib
matplotlib.rcParams['font.family'] = 'SimHei'

L = 1.0
T = 100.0
rho = 0.01
c = np.sqrt(T / rho)
nx = 101
dx = L / (nx - 1)
dt = 0.5 * dx / c
nt = 2000
x = np.linspace(0, L, nx)

u = np.zeros((3, nx))
u[0, :] = np.sin(np.pi * x / L) * 0.1

for n in range(1, nt):
    for i in range(1, nx - 1):
        u[2, i] = (c * dt / dx) ** 2 * (
            u[1, i + 1] - 2 * u[1, i] + u[1, i - 1]
        ) + 2 * u[1, i] - u[0, i]
    u[2, 0] = 0.0
    u[2, -1] = 0.0
    u[0, :] = u[1, :]
    u[1, :] = u[2, :]

fig, ax = plt.subplots(figsize=(10, 3))
line, = ax.plot(x, u[2, :], 'b-', lw=2)
ax.set_xlim(0, L); ax.set_ylim(-0.15, 0.15)
ax.set_xlabel('x'); ax.set_ylabel('y(x,t)')
ax.set_title('两端固定弦的振动')

def update(n):
    for i in range(1, nx - 1):
        u[2, i] = (c * dt / dx) ** 2 * (
            u[1, i + 1] - 2 * u[1, i] + u[1, i - 1]
        ) + 2 * u[1, i] - u[0, i]
    u[2, 0] = 0.0; u[2, -1] = 0.0
    u[0, :] = u[1, :]; u[1, :] = u[2, :]
    line.set_ydata(u[2, :])
    return line,

ani = FuncAnimation(fig, update, frames=500, interval=20, blit=True)
plt.show()
\end{lstlisting}
\end{mycode}

\begin{mycode}{绘制弦的前五阶振型}
\begin{lstlisting}[language=Python]
import numpy as np
import matplotlib.pyplot as plt

L = 1.0
x = np.linspace(0, L, 200)

fig, ax = plt.subplots(figsize=(8, 6))
colors = ['blue', 'red', 'green', 'orange', 'purple']

for n in range(1, 6):
    Y_n = np.sin(n * np.pi * x / L)
    ax.plot(x, Y_n + 0.2 * (n - 1), color=colors[n-1], lw=2,
            label=f'n={n}')

ax.set_xlabel('x'); ax.set_ylabel('Y_n(x)')
ax.set_title('两端固定弦的前五阶固有振型')
ax.legend(); ax.grid(True, alpha=0.3)
plt.tight_layout(); plt.show()
\end{lstlisting}
\end{mycode}

\begin{mycode}{不同边界条件下弦的固有频率对比}
\begin{lstlisting}[language=Python]
import numpy as np

L = 1.0; T = 100.0; rho = 0.01; c = np.sqrt(T / rho)

print("边界条件        固有频率 (Hz)")
print("-" * 40)

n_values = np.arange(1, 6)
omega_ff = n_values * np.pi * c / L
print("两端固定:      ", np.round(omega_ff / (2 * np.pi), 2))

omega_cf = (2 * n_values - 1) * np.pi * c / (2 * L)
print("一端固定一端自由:", np.round(omega_cf / (2 * np.pi), 2))

omega_elas = n_values * np.pi * c / L
print("两端自由(弹性): ", np.round(omega_elas / (2 * np.pi), 2))
\end{lstlisting}
\end{mycode}

--------------------------------------------------------------------
\section{小结}

本章从离散系统过渡到连续系统，建立了弦和杆的一维波动方程。核心要点：

\begin{enumerate}
    \item 连续系统由偏微分方程描述，拥有无穷多自由度;边界条件内在于解中;
    \item 一维波动方程 $\partial^2\phi/\partial t^2=c^2\partial^2\phi/\partial x^2$ 是弦、杆纵振、杆扭振的统一数学形式;
    \item 分离变量法将 PDE 分解为空间 ODE 和时间 ODE，边界条件确定固有频率和振型;
    \item 振型满足正交性，这使得模态叠加法可以将 PDE 解耦为无穷多个单自由度方程;
    \item 数值方法（如有限差分法）是求解复杂边界和初始条件下波动问题的有力工具。
\end{enumerate}

